\begin{frame}{Quem sou eu}
  \begin{columns}
    \column{0.6\textwidth}
    {\Large\textbf{Silas Henrique}}
    \vspace{0.5cm}
    \begin{itemize}
      \item Técnico em Eletromecânica pelo IFBA
      \item Graduando em Engenharia de Controle e Automação (8º semestre)
      \item Ex-diretor de Marketing da OPTIMUS Jr.
      \item Pesquisador do PRH 41/UFBA
    \end{itemize}
    \vspace{0.5cm}
    \href{https://silas.eng.br/}{\textbf{🌐 silas.eng.br}}
    \column{0.4\textwidth}
    \centering
    \begin{tikzpicture}
      \clip (0,0) circle (0.5\textwidth);
      \node at (0,0) {\includegraphics[width=\textwidth]{me.png}};
    \end{tikzpicture}
  \end{columns}
\end{frame}

\begin{frame}{Antes de começarmos}
  \begin{block}{Ritmo do curso}
    \begin{itemize}
      \item O curso é \textbf{acelerado}
      \item Eu costumo \textbf{falar rápido}
      \item Se eu passar rápido demais, \textbf{me avisem}
    \end{itemize}
  \end{block}
  \begin{alertblock}{Sobre o rigor teórico}
    \begin{itemize}
      \item O foco aqui é \textbf{entender e usar os conceitos}
      \item Não vamos entrar em todo o rigor da teoria
      \item Em disciplinas como \textbf{Física III} vocês verão definições mais formais
    \end{itemize}
  \end{alertblock}
\end{frame}
